\section{Empirical Falsification Boundary}

\subsection{Why an empirical boundary is needed}


The native finite Thalean relational-saturation theorem is now an
internal mathematical result.

The present paper does not claim empirical confirmation or physical
realization of that theorem.

The wider research program nevertheless contains a useful
methodological precedent: a native finite registration was converted
into a bounded scalar detector and subjected to a frozen
cross-instrument replication protocol.

That precedent shows how a future physical saturation test should be
designed.

It does not provide empirical evidence for the finite saturation
theorem itself.

\subsection{The frozen G60 C3 detector}

The WMAP replication \citep{cave2026wmap} reused the native
\(C_3\) registration inherited from the earlier Planck experiment.

For each accepted spherical placement and retained multipole band, the
detector evaluates three registered phase states under the native
\(C_3\) action.

Their harmonic response energies define a local three-phase profile.

That profile is mapped to a bounded scalar
\[
\Xi.
\]

The real-data pipeline is schematically
\[
T_{\ell_{\max}}
\longrightarrow
\nabla_{S^2}T_{\ell_{\max}}
\longrightarrow
\text{registered }C_3\text{ edge responses}
\longrightarrow
\text{three-phase energy profile}
\longrightarrow
\Xi.
\]

The exact graph registration, phase order, orientation, spherical
sampling geometry, edge-field construction, energy definition, and
\(\Xi\) evaluator were frozen before WMAP unblinding.

\subsection{Preregistered replication rule}

The frozen primary target was
\[
\ell_{\max}=64,
\qquad
\text{mean }\Xi,
\qquad
\text{upper tail}.
\]

Both WMAP V and W products were required independently to satisfy
\[
p_{\mathrm{two}}\le0.05
\]
against both of their frozen null families:

\begin{enumerate}[label=(\roman*)]

\item Gaussian nulls;

\item amplitude-preserving phase-surrogate nulls.

\end{enumerate}

Thus all four required product-null comparisons had to pass.

No combined probability was permitted.

No neighboring scale, alternate reduction, secondary product, or
post-hoc statistic could rescue a failed primary comparison.

\subsection{Observed result}

The four frozen primary comparisons were

\[
p_{\mathrm{V,Gaussian}}
=
0.0239760,
\]

\[
p_{\mathrm{V,phase}}
=
0.00799201,
\]

\[
p_{\mathrm{W,Gaussian}}
=
0.0539461,
\]

and

\[
p_{\mathrm{W,phase}}
=
0.0339660.
\]

All four real WMAP statistics lay in the preregistered upper-tail
direction.

Three of the four comparisons crossed
\[
p\le0.05.
\]

The WMAP W Gaussian comparison did not.

Therefore the frozen decision rule yields
\[
\boxed{\text{NO-GO REPLICATION}}.
\]

No post-hoc rescue is applied.

\subsection{What the WMAP result does and does not show}


The result preserves one descriptive observation:

\begin{quote}
The predicted upper-tail direction appeared in all four required WMAP
comparisons.
\end{quote}

But the bounded cross-instrument replication claim was not earned.

The finite mathematical construction is unchanged by that outcome.

The WMAP paper therefore provides a methodological lesson, not
empirical support for native relational saturation, quadratic contrast
concentration, or any physical saturation bridge.

\subsection{No retroactive saturation detector}

The scalar
\[
\Xi
\]
was not preregistered as a relational-saturation observable.

It measured a frozen \(C_3\) phase-structure response.

Therefore Paper 14 does not reinterpret
\[
\Xi
\]
after the fact as:

\begin{itemize}

\item a saturation measure;

\item a gravity measure;

\item a black-hole proxy;

\item an information-density measure;

\item or a contrast-concentration statistic.

\end{itemize}

Doing so would define a new experiment.

\subsection{Requirements for a future saturation test}


A future empirical test must first declare which finite theorem-level
operation is being mapped into physics.

At minimum it should freeze:

\begin{enumerate}[label=(\roman*)]

\item the physical or complete mathematical state family;

\item the declared projection or registration;

\item the proposed physical analogue of native selection, if any;

\item the proposed contrast-concentrating operation, if any;

\item the aggregate observable;

\item the registered contrast observable;

\item the expected relation among those quantities;

\item the null model or comparison family;

\item the predicted direction;

\item the decision threshold;

\item and the rule for classifying the outcome.

\end{enumerate}

Native relational loading and quadratic contrast concentration may not
be treated as one undifferentiated loading variable.

A proposed physical cycle containing both must specify their coupling
before observing the data.

\subsection{A minimal structural observable}


Suppose a physical bridge supplies a declared register
\[
h(\lambda)
\]
indexed by a physical control parameter
\[
\lambda.
\]

Define
\[
C(\lambda)=\|P_0h(\lambda)\|_2
\]
and
\[
\zeta(\lambda)
=
\frac{\|P_\perp h(\lambda)\|_2^2}
{\|h(\lambda)\|_2^2}.
\]

A direct test of a \emph{contrast-concentrating} physical regime might
predict that
\[
C(\lambda)
\]
remains large or increases while
\[
\zeta(\lambda)
\]
decreases.

That is not the theorem-level signature of native relational loading.

Under the canonical finite completion register, native loading
strictly increases \(\zeta\).

Therefore any physical experiment must state whether it is testing:

\[
\text{native future selection},
\]
\[
\text{contrast concentration},
\]
or
\[
\text{an explicitly proposed composition of the two}.
\]

The existence of an observed pattern would test the declared bridge.

It would not, by itself, establish its physical cause.

\subsection{Falsification logic}


A future physical saturation hypothesis should count against itself if
one or more of the following occurs under a valid frozen protocol:

\begin{enumerate}[label=(\roman*)]

\item the declared native-selection analogue fails to produce its
predicted change in future or state multiplicity;

\item the declared contrast-concentrating analogue fails to produce
its predicted contrast behavior;

\item a proposed coupling between the two is not observed;

\item the observed relation is reproduced equally by declared null
structures lacking the relevant Thalean organization;

\item the signal depends on post-hoc changes to registration, loading,
scale, coupling law, or threshold;

\item or the proposed effect fails independent replication.

\end{enumerate}

The purpose is to keep the physical extension falsifiable rather than
interpretively elastic.

\subsection{Boundary}


The exact empirical status carried into Paper 14 is:

\[
\text{native finite detector methodology}
\quad
\text{demonstrated},
\]

\[
\text{WMAP directional concordance}
\quad
\text{observed},
\]

\[
\text{WMAP preregistered replication}
\quad
\text{NO-GO},
\]

\[
\text{physical relational saturation}
\quad
\text{untested}.
\]

No empirical result currently promotes the native finite
relational-saturation theorem to a physical claim.
